% !TEX program = xelatex
% v3 中文镜像稿（架构：plan/architecture.md；对应 en/main.tex 逐节镜像）
% 编译：xelatex main -> bibtex main -> xelatex main -> xelatex main
\documentclass[UTF8,a4paper,10pt,twocolumn]{ctexart}

\usepackage{amsmath,amssymb}
\usepackage{booktabs}
\usepackage{graphicx}
\usepackage{subcaption}
\usepackage{multirow}
\usepackage{array}
\usepackage{url}
\usepackage[hidelinks]{hyperref}
\usepackage{enumitem}
\usepackage{algorithm}
\usepackage{algpseudocode}

\usepackage{geometry}
\geometry{a4paper,left=1.6cm,right=1.6cm,top=2.0cm,bottom=2.0cm,columnsep=0.65cm}
\setlength{\parindent}{2em}
\linespread{1.06}

\renewcommand{\figurename}{图}
\graphicspath{{../figures/}}
\renewcommand{\tablename}{表}
\captionsetup{font=small,labelfont=bf,labelsep=quad}

\newcommand{\dblfig}[3]{%
  \begin{figure*}[!t]\centering
    \includegraphics[width=\textwidth]{#1}%
    \caption{#2}\label{#3}
  \end{figure*}}
\newcommand{\onefig}[4]{%
  \begin{figure}[!t]\centering
    \includegraphics[width=\columnwidth]{#1}%
    \caption{#2}\label{#3}\vspace{#4}
  \end{figure}}

\title{\vspace{-1.2em}\heiti\zihao{3} 首轮是一个双向窗口：QC-MDPC\\比特翻转译码中的阈值时序与条件失败}

\author{\zihao{5} [作者姓名待补充]\\[0.2em]
\zihao{6}\mdseries [单位与通讯邮箱待补充]}

\date{}

\begin{document}

\maketitle

\begin{abstract}
BIKE 等 QC-MDPC 方案的安全性论证依赖一个任何直接模拟在可行样本量下都无法解析的失败率目标（约 $2^{-128}$），也依赖译码行为在无法逐一检查的密钥之间保持稳定。本文考察 BF 译码器的阈值调度能在多大程度上控制条件失败。我们把两条调度施加于完全相同的密钥、错误与综合征（配对实例），并在全部推断中重抽密钥而非错误。CLASSIC BF 译码器的首轮是一个陡峭且不对称的双向窗口：把其阈值提高 1 至 3，三个正式条件上的条件失败率下降 2.3--5.5 倍；降低 1，恶化约八倍；降低 2，恶化超过两个数量级。偏移的轮次与持续方式和幅度同等重要：同一偏移只作用于自身那一轮时有益，持续贯穿后续轮次时则使失败率升至 1.2--6.1\%，即便第 1 轮未被触及。首轮记录解释了这一不对称：错误的首轮翻转下降得远多于正确翻转，即干预剔除的是错误决策，而非加速纠错。效应无需重调即迁移到第二个工作点与更大的码，并在所测的捕获倾向构造处止步。在生产块尺寸上，错误重量一经抬高，窗口即复现；在生产重量本身，32{,}000 次译码的零失败只给出 rule-of-three 上界，另有决策余量与收敛代价可测。一个首轮代理模型未能预测另一套调度的效应；我们把处理后条件化认定为失败模式。
\end{abstract}

\vspace{0.6em}

\noindent\textbf{关键词：}后量子密码；QC-MDPC 码；比特翻转译码；解码失败率；首轮阈值；阈值时序；密钥条件风险。

\vspace{0.8em}

\section{引言}\label{sec:intro}

BIKE 等 QC-MDPC 方案建立在迭代式比特翻转（BF）译码之上，其罕见的失败塑造着它们的安全论证~\cite{misoczki2013mdpc,bike2024spec,alagic2025nist}。Level-1 的目标是解码失败率（DFR）约 $2^{-128}$，远低于直接仿真所能触及的水平~\cite{sendrier2019dfr,vasseur2021thesis}。失败并不均匀铺开：它们聚集在特定密钥上~\cite{guo2016key,wang2023weak,nosouhi2023weakkey}，而能检测或计时这些失败的攻击者可恢复密钥~\cite{eaton2018timing}。译码器逐轮的阈值构成一套调度，调整它无须触碰密钥生成。因此我们检验一个带条件的问题：对给定密钥，刻意改变调度是否会改变失败风险。

在刻画 Vasseur 参考实现中的 CLASSIC BF 译码器~\cite{vasseur2021thesis}（记作 C0）时，我们观察到它最早的那次决策异常敏感。首轮作用于未经触碰的综合征，因此首轮的任何失误都会进入此后每一轮的状态。规范已经指出首轮是失败的主要来源~\cite{bike2024spec}；尚未测量的是首轮阈值如何改变条件失败，以及该效应如何依赖轮次与持续方式。把首轮阈值加 2，条件失败率在随机条件下降为原来的 1/4.15，在对抗探针下降为 1/5.47；把阈值减 1，失败率反而上升约八倍。两单位偏移若从第 2 轮起持续施加则具破坏性，只限于一轮时却有益。

我们围绕配对实例重建实验。一个实例是一把密钥 $H$、一个错误 $e$ 及其综合征 $s=He^{\top}$；每个调度变体都译码同一个三元组。FRO-$k$ 只在第 1 轮把阈值抬高 $k$，FRONEG 则把阈值减 1。比较用风险比（RR）汇总，即基线失败率除以该变体的失败率，取值大于 1 有利于该变体。由于失败聚集在密钥上，区间与检验重抽样密钥而非错误。主要效应估计位于失败频繁到足以计数的区间；零失败单元以 rule-of-three 上界报告，我们对生产尾部不作主张。

本研究有四项贡献和一个设计好的负结果。第一，首轮是一扇双向窗口。保守偏移 $+1$ 至 $+3$ 使失败率下降 2.3--5.5 倍，而宽松侧在两个计数内崩塌（第~\ref{sec:formal} 节）。第二，同一偏移的符号随施加轮次与持续方式改变。$+2$ 只在第 2 轮施加一次时带来小收益，从第 2 轮起持续施加则给出 RR 0.0168（第~\ref{sec:temporal} 节）。第三，机制是不对称的。FRO-2 把错误的首轮翻转降低约十倍，而正确翻转降得少得多（第~\ref{sec:mechanism} 节）。第四，只要把错误重量抬进可测区间，窗口就在生产块尺寸上复现；在生产重量本身，只有决策余量与收敛代价可测（第~\ref{sec:production} 节）。这个设计好的负结果收束清单：一个首轮代理模型未能预测另一套调度的效应（第~\ref{sec:counterfactual} 节）。

第~\ref{sec:background} 节搭建 QC-MDPC BF 译码、相关工作与首轮窗口的一阶模型。第~\ref{sec:method} 节给出配对设计，第~\ref{sec:results} 节给出结果。第~\ref{sec:discussion} 节解读结果并陈述局限；第~\ref{sec:conclusion} 节收束。
\section{背景与相关工作}\label{sec:background}

本节固定码、译码器、我们所对照的已有工作，以及首轮窗口的一阶模型。

\subsection{QC-MDPC 码与 BF 译码器}

准循环中密度奇偶校验（QC-MDPC）码是校验矩阵块循环的二进制线性码~\cite{misoczki2013mdpc}。$r$ 为块长，$n=2r$ 为码长。校验矩阵为 $H=[H_0\mid H_1]$，每个循环块的行重为 $d$；完整矩阵的行重为 $2d$，列重为 $d$。$t$ 为一次译码尝试注入的错误数。主工作点为 $(r,d,t)=(4801,45,92)$（第~\ref{sec:method} 节）。对校验矩阵 $H$ 与错误向量 $e$，综合征为 $s=He^{\top}$。最终综合征为零则译码成功，否则失败。

BF 译码迭代地纠正错误。一个位置的决策变量是它的未满足校验计数（UPC），即涉及它的未满足方程数。每一轮译码器翻转所有 UPC 达到阈值的位置，重算综合征，如此重复，直到综合征消失或迭代预算用尽。算法~\ref{alg:bf} 给出该过程。

\begin{algorithm}[t]
\caption{普通 BF 译码}
\label{alg:bf}
\begin{algorithmic}[1]
\Require $H$、输入综合征 $s_{\mathrm{in}}=He^{\top}$、预算 $I$
\Ensure 错误估计 $\hat e$，或失败
\State $\hat e \gets \mathbf{0}$；$s_{\mathrm{res}} \gets s_{\mathrm{in}}$
\For{$i=1,\ldots,I$}
  \State 对每个位置 $j$ 计算 $U_j=\mathrm{UPC}(H,s_{\mathrm{res}},\hat e,j)$
  \State 由当前综合征重量 $|s_{\mathrm{res}}|$ 计算 $\tau_i$
  \State 翻转所有满足 $U_j\ge\tau_i$ 的 $j$
  \State 更新 $\hat e$，并令 $s_{\mathrm{res}}\gets s_{\mathrm{in}}\oplus H\hat e^{\top}$
  \If{$s_{\mathrm{res}}=\mathbf{0}$}
    \State \textbf{return} $\hat e$
  \EndIf
\EndFor
\State \textbf{return} 失败
\end{algorithmic}
\end{algorithm}

基线是 Vasseur 参考实现中的 CLASSIC BF 译码器~\cite{vasseur2021thesis}，记作 C0。其阈值由 \path{compute_threshold} 的二项模型给出，向上取整并裁剪到 $[(d+1)/2,d]$，每轮根据当前综合征重量重算。主实验使用 20 轮预算与静态距离语义。静态运行使用固定的初始 $t$；偏移加到裁剪后的基础阈值上，再裁剪回同一区间。随代码归档的 \path{reference/FROZEN_REFERENCE_SEMANTICS.md} 给出精确公式与来源身份。零综合征或用尽预算时停止。研究载体在一轮没有翻转时继续到预算。C 参考实现在这种轮次提前退出；对 C0 两者结果等价，因为状态不再变化，但一次性偏移会在下一轮被移除，因此继续规则对 FRO-2 有影响。本文全程使用继续规则。另有一套单独编写的动态距离参考实现，它使用当前到真实错误的距离（部署者不知道这一信息），只作为仿真特权的方向性核对。

它有名字的后继修改调度或携带额外状态：Black-Gray-Flip（BGF）、Backflip，以及规范中的译码器，我们称之为 BIKE-Flip~\cite{drucker2020shades,vasseur2021thesis,bike2024spec}；Bartz 与 Liva 比较了更早的 MDPC-McEliece 译码方案~\cite{bartz2019decoding}。我们把它们当作彼此不同的译码器。BIKE-Flip 在每一轮都施加下限保护，$T=\max(f_t(S),T_i(S_0))$~\cite{bike2024spec}，其中 $S$ 为当前综合征重量，$f_t(S)$ 为其仿射函数，$T_i(S_0)$ 为由初始综合征重量 $S_0$ 计算的第 $i$ 轮下限。每个下限等于未加裕度的水平加上裕度 $\Delta$，因此裕度持续存在于每一轮，而不只出现在首轮。Level-1 规范取 $\Delta=3$，预算为七次迭代。

规范附录 A.2.4 以平均综合征重量处理论阈值的切线推导仿射常数；我们用规范为 Level 1、3、5 印出的常数核对了实现。已发表的机理解释把残余失败放在第 2、3 轮：综合征重量已经下降，阈值跌破少数正确位置的计数，译码器便把正确位置当作错误翻转~\cite{sendrier2024seminar,wang2023weak}。

FRO-2 走的是另一条路线：它只在首轮把阈值加 $+2$ 一次，后续轮保持普通阈值。

\subsection{失败以密钥与调度为条件}

失败不仅以信道为条件，也以密钥为条件。Guo 等表明触发失败的错误图样与秘密矩阵相关，因此观测到的失败泄露密钥信息~\cite{guo2016key}；Wang 等把同一相关性变成弱密钥构造~\cite{wang2023weak}；Eaton 等把计时观测转成密钥恢复攻击~\cite{eaton2018timing}；Gellersen 等扩展了侧信道路线~\cite{gellersen2025}。失败率偏高的弱密钥类也已被构造、枚举和筛选~\cite{vasseur2021thesis,nosouhi2023weakkey,matthews2025combinatorial}，弱密钥攻击也围绕 BGF 译码得到优化~\cite{liu2025bgf}。调度是第二个条件变量，因为一轮翻转的位置集合塑造进入下一轮的状态。我们的设计因此扰动调度，并把密钥作为推断单元：所有变体译码相同的 $(H,e,s)$ 实例，不确定度来自重抽密钥而非错误（第~\ref{sec:method} 节）。

\subsection{相关工作}

失败建模沿两条路线展开。第一条从译码器状态解释失败。Kim 与 Park 拟合逐实例逻辑预测器~\cite{kim2021logistic}；Paiva 与 Terada 用首轮残余错误作为代理~\cite{paiva2022pickyfix}；Santini 等与 Annechini 等分析正则码的计数分布~\cite{santini2019conservative,santini2020analysis,annechini2024estimation}。第二条估计仿真无法触及的失败率。Sendrier 与 Vasseur 把 BF 建模为马尔可夫链~\cite{sendrier2019dfr,sendrier2020lowdfr}；Tillich 给出失败概率界~\cite{tillich2018dfp}；Vasseur 加入带置信区间的外推~\cite{vasseur2021thesis}；Arpin 等预测密钥相关错误平层~\cite{arpin2025errorfloor}；Annechini 等给出闭式估计，并明确警告基于代理的外推是危险的~\cite{annechini2025closedform}。

第三条路线直接调阈值。Kichna 等分析不同阈值公式如何改变迭代数与 DFR（跨错误重量与安全级别）~\cite{kichna2024thresholds}，随后在我们这一块长上用粒子群搜索把推荐错误重量下的平均迭代从 5.36 降到 3.01~\cite{kichna2025pso}。这些研究比较整条调度的选择；它们没有隔离一次性首轮位移，也没有在相同 $(H,e,s)$ 实例上配对调度。近期工作改变译码器本身：BF-Max~\cite{baldelli2025bfmax}、原地随机翻转~\cite{baldi2021inplace}、近码字感知 BF~\cite{baldelli2026ncw}、分层 BF~\cite{cai2025layered}，以及后续改进~\cite{santini2026flip}。这些译码器是量级与代价的参照，不是我们声称支配的对象。这些路线共同引出本文的问题：首轮调度的微小改变会改变条件失败吗？效应是否依赖施加轮次与持续方式？我们在配对实例上检验这一问题；第~\ref{sec:counterfactual} 节再给出一个次级诊断，检验同一批首轮统计量能否预测另一套调度。
\subsection{首轮窗口的一阶模型}\label{sec:model}

模型刻意保持最小。它描述首轮阈值的轻微抬高如何重塑该轮的翻转；后续轮不在其范围内。

固定一个实例，即一把密钥 $H$、一个错误 $e$ 及其综合征 $s$。首轮中每个位置 $j$ 携带计数 $U_j$，即它的 UPC。译码器翻转计数达到首轮阈值 $\tau_1$ 的位置，其余不翻。

错误位置与干净位置在译码器眼中并不相同。错误支持集内的位置把自身错误带入每个关联校验，因此多数校验不满足；干净位置的校验只因其它错误而不满足。在首轮 UPC 分布所用的独立性近似下~\cite{annechini2024estimation}，两个总体服从不同的二项分布，每校验概率为 $p_1>p_0$，式~\eqref{eq:upc_tails}：
\begin{equation}\label{eq:upc_tails}
\begin{aligned}
U_{i,j} &\sim \operatorname{Bin}(d,p_i), \qquad i\in\{1,0\},\; p_1>p_0,\\
q_i(\tau) &= \Pr[U_{i,j}\ge\tau]
= \textstyle\sum_{u=\tau}^{d}\binom{d}{u}p_i^{u}(1-p_i)^{d-u}.
\end{aligned}
\end{equation}
其中 $i=1$ 标记错误位置，$i=0$ 标记干净位置。尾质量 $q_i(\tau)$ 是第 $i$ 类位置在阈值 $\tau$ 处翻转的概率，因此 $q_1$ 支配正确翻转，$q_0$ 支配假翻转。这些概率是模型参数而非拟合量；下面的乘子由归档的首轮普查计算，模型只作定性使用。

扰动是一次性抬高 $\tau_1\to\tau_1+2$。抬高通过幸存翻转比例作用于每个总体，式~\eqref{eq:multipliers}：
\begin{equation}\label{eq:multipliers}
A_1=\frac{q_1(\tau_1+2)}{q_1(\tau_1)},\qquad
A_0=\frac{q_0(\tau_1+2)}{q_0(\tau_1)} .
\end{equation}
模型预测 $A_0\ll A_1$：$\tau_1$ 处的干净尾是更薄的外尾，两个额外计数因此不成比例地切进它。在留出点 $P_4=(3607,39,76)$，归档剖面给出 $A_1 = 0.355$、$A_0 = 0.0892$。两者都小于一，所以抬高同时削薄两个总体；远小的 $A_0$ 把削薄主要放在假翻转上，即有利方向。被推迟的正确翻转仍是代价。

模型给出三个可检验推论，组织后续实验。
\begin{enumerate}[label=(\roman*),leftmargin=*,itemsep=1pt,topsep=2pt]
\item 首轮普查应直接显示不对称：假翻转的相对下降应远大于正确翻转（第~\ref{sec:mechanism} 节）。
\item 净失败变化没有局部固定的符号，因为被推迟的正确翻转是反向代价；必须由配对窗口裁决（第~\ref{sec:results} 节）。
\item 模型只约束一轮，因此持续偏移是另一项实验（第~\ref{sec:temporal} 节）。
\end{enumerate}
模型对净失败变化的符号和首轮之后的轮次保持沉默。推论 (ii) 与 (iii) 因此是实验而非推论。
\section{研究设计}\label{sec:method}

第~\ref{sec:model} 节的模型没有规定首轮窗口的符号；本节固定测量它的实验方案。基线是 C0，干预是 FRO-2，两者均在第~\ref{sec:intro} 节定义。

对每个条件中的每个变体，我们以失败译码次数除以尝试次数估计 DFR。区间不确定度来自对密钥的 $95\%$ 自助法，而不是对错误抽取的自助法，因此区间反映密钥之间的变动。零失败条目绝不被读成零风险；在 $N$ 次尝试上改用 rule-of-three 上界 $3/N$。每次比较都归结为风险比，即 C0 的 DFR 除以该变体的 DFR，取值大于一有利于该变体。各变体只在阈值调度上不同；迭代预算及其余设置固定不变。

区间与检验口径在运行之前固定。区间是跨密钥的百分位聚类自助区间：每个对比 10{,}000 次重抽样，密钥有放回抽取，被抽中密钥的全部错误抽取一并保留，自助种子固定为 20260905。符号检验统计该变体失败数少于 C0 的密钥数，丢弃持平密钥，并给出精确双侧二项概率。Holm 校正覆盖主对比的三个正式条件；时序网格、迁移检验与译码器对比以边际区间报告。所有计算都可由归档的逐钥计数复现。

本研究使用八个参数点；生产块报告两个错误重量。主工作点是 $P_1'=(4801,45,92)$；撇号标记筛选点 $P_1=(4801,45,84)$ 的可测变体，候选调度在正式测量之前于该筛选点比较。迁移点是 $P_2=(3001,37,66)$ 与 $P_3=(4001,41,82)$，$P_4=(3607,39,76)$ 留作规则检验。系列中最大的块是 $(6007,56,100)$。在生产块 $(12323,71)$ 上，本研究测量抬高重量 $(t^{\ast}=172)$ 与生产重量 134；CPA-128 预设为 $(10163,71,140)$。

实验使用三种密钥/错误采样条件。RAND 随机抽取密钥，并在工作错误重量 $t$ 上均匀随机抽取错误。TYPE-I 约束密钥的某个校验块含 $12$ 个连续的 1，错误仍按 RAND 抽取；它用作对抗性探针，而不是总体估计。NCW-shallow 使用经过筛选、允许近码字构造的密钥池，并构造交叠量 $\kappa=12$ 的错误，低于块行重 $d=45$。浅交叠条件是一种压力构造，并不证明捕获不可能；另一族跨越 $d$ 的压力测试见第~\ref{sec:boundary} 节。图~\ref{fig:design} 展示配对与运行流。

实验分阶段推进，从标定一直到在新密钥上做一次确认。我们先在主密钥族上标定，那里失败可测，使每个被比较的差异都升到零失败下限之上。预指定之前，我们以最差组上界为准则筛选了九个候选调度；该上界读取结果最弱的条件，调度因此无法把某个条件的危害藏在另一个条件的收益背后。预指定门槛要求最差组的 $95\%$ 上界至少改善 $0.5$ 个 $\log_{10}$ 单位（约 3.2 倍），此外主条件必须保持非劣。门槛随后在纯新子集上评估一次：212 把密钥，每把 3000 次抽取，每变体 636{,}000 次。

这些密钥从未进入筛选，因此确认独立于每一项选择决定；该子集仍属于 278 把密钥的正式池，所以这一独立性是针对筛选过程而言，而不是第二次统计独立的估计。四个冻结主种子分别管辖主流程、自助、探索与确认四条流，因此每次运行都可精确复现。

后文报告的每个数字都可由冻结导出复现。同一项目内独立编写的第二套实现，在 20{,}000 次动态距离译码中记录到 11 次失败，C 参考则是 200{,}000 次中 136 次（$z=-0.68$），并复现了主对比的方向。这一交叉核对只能限定编码错误，因为两套实现共享同样的设计假设。另有一套静态载体的 C 实现，基于参考代码并采用继续规则与相同首轮偏移，在三组正式池上逐钥复现了 Python 载体的失败计数：RAND 282/68、TYPE-I 372/68、NCW-shallow 141/31。由于它读取同一批导出实例，这是对主估计量的实现无关核对，而不是第二次总体估计。精确恢复审计覆盖全部 278 把正式 $P_1'$ 密钥，每变体 834{,}000 次译码。它发现 C0 有 795 次精确恢复失败，FRO-2 有 167 次，与零综合征总数吻合，两个变体都没有观察到错纠。主实验在静态距离语义下译码，C 参考使用动态距离，因此两者的失败率从不直接比较。探索性条目与归档原型保持标注，从不承载主结论。

\begin{figure*}[!t]
\centering
\includegraphics[width=0.85\textwidth]{fig1_study_design.pdf}
\caption{研究设计。(a)~两条调度译码同一把密钥、同一个错误与同一个综合征。(b)~FRO-2 只在首轮加二，从第 2 轮起两者都使用 C0 调度。(c)~结果按密钥配对，密钥是重抽样单元。}
\label{fig:design}
\end{figure*}
\section{结果}\label{sec:results}

本节按固定顺序报告证据：窗口、它的时序、机制、迁移、边界，以及现代译码器参照。

\subsection{首轮是一个双向窗口}\label{sec:formal}

每个实例由 C0 与只在首轮阈值加 $+2$ 的 FRO-2 分别译码。两个变体看到同一把密钥、同一个错误与同一个综合征。对比在 $P_1'=(4801,45,92)$ 上进行，覆盖 RAND、TYPE-I 与 NCW-shallow。表~\ref{tab:main} 给出计数与检验；图~\ref{fig:main} 给出绝对率、区间与逐钥分布。

FRO-2 在三种条件下都降低失败率（表~\ref{tab:main}）。RR 为 4.15 [3.29, 5.38]、5.47 [4.25, 7.28] 与 4.55 [3.21, 6.82]。每个区间都排除 1，密钥聚类符号检验在 Holm 校正后也一致。逐钥看，FRO-2 在 278 个密钥-条件情形中改善 233 个、变差 18 个、持平 27 个。预指定门槛要求最差条件的 $95\%$ DFR 上界至少改善 $0.5$ 个 $\log_{10}$ 单位（约 3.2 倍）。全池改善 0.679 个 $\log_{10}$（约 4.8 倍），纯新子集改善 0.687 个 $\log_{10}$，因此门槛通过。

窗口两侧并不对等。FRONEG 在首轮减 1：它使 92 把 RAND 密钥全部变差，失败数从 282 升到 2{,}235，RR 为 0.126 [0.111, 0.141]，即上升 7.93 倍。向下走一步的代价因此超过向上走两步的收益。

完整局部扫描分辨出两侧（图~\ref{fig:window}）。偏移 $+1$ 至 $+3$ 使失败率下降约 2.3 至 5.5 倍，最大值落在 $+2$ 或附近。宽松侧在两个计数内崩塌：降低 1 给出 RR 0.10--0.13，降低 2 给出 0.0014--0.0030。在 RAND 下，两计数下调把失败数放大 365 倍，从 276{,}000 次中的 282 次到 102{,}980 次。窗口是双向且陡峭不对称的。直白地说，保守偏移换来 4--5 倍的下降，而向下走一步的代价超过这一收益。

\begin{table*}[!t]
\centering
\small
\caption{$P_1'=(4801,45,92)$ 处的主对比。单元格给出失败数/译码数（DFR）。RR 为 C0 除以 FRO-2，带 95\% 密钥聚类自助区间；Holm $p$ 为跨三条件的密钥聚类符号检验。}
\label{tab:main}
\begin{tabular}{l r r r r r}
\toprule
条件 & 密钥数 & C0 失败数/译码数 (DFR) & FRO-2 失败数/译码数 (DFR) & RR [95\% CI] & Holm $p$ \\
\midrule
RAND & 92 & 282/276{,}000 ($1.02\times10^{-3}$) & 68/276{,}000 ($2.46\times10^{-4}$) & 4.15 [3.29, 5.38] & $4.88\times10^{-17}$ \\
TYPE-I & 94 & 372/282{,}000 ($1.32\times10^{-3}$) & 68/282{,}000 ($2.41\times10^{-4}$) & 5.47 [4.25, 7.28] & $3.05\times10^{-22}$ \\
NCW-shallow & 92 & 141/276{,}000 ($5.11\times10^{-4}$) & 31/276{,}000 ($1.12\times10^{-4}$) & 4.55 [3.21, 6.82] & $7.66\times10^{-12}$ \\
\bottomrule
\end{tabular}
\end{table*}

\dblfig{fig2_main_results.pdf}{FRO-2 在三种条件下都降低观测失败率。(a)~每 10{,}000 次译码的绝对失败数。(b)~C0/FRO-2 风险比与 95\% 密钥聚类自助区间，虚线为 1。(c--e)~RAND、TYPE-I 与 NCW-shallow 下逐钥避免失败数的分布；标注给出改善、变差与持平的密钥数。每把密钥每条调度 3{,}000 次译码。}{fig:main}

\begin{figure*}[!t]
\centering
\includegraphics[width=0.88\textwidth]{fig6_threshold_window.pdf}
\caption{$P_1'$ 处首轮是一扇双向窗口。(a)~六个偏移的绝对 DFR（对数轴）。(b--d)~正偏移 $+1$、$+2$、$+3$ 按条件给出的风险比，统一从 0 起的线性轴，带 95\% 密钥聚类区间；浅色带只标出 $+2$，不表示统计优越性。负偏移的绝对率见 (a)。}
\label{fig:window}
\end{figure*}

\subsection{偏移的轮次与持续方式}\label{sec:temporal}

轮次对照改变偏移的施加轮次与持续方式。R2one 只在第 2 轮加 $+2$；R2plus 从第 2 轮起持续加 $+2$。同样的幅度在首轮给出 RR 4.15，在第 2 轮给出 1.21，持续到后续轮则给出 0.0168。

在 RAND 下，R2one 给出 1.21 [1.07, 1.36]，44 把密钥改善、25 把变差、23 把持平。R2plus 给出 0.0168 [0.0149, 0.0187]：失败数从 282 升到 16{,}785，DFR 升到 $6.08\times10^{-2}$，92 把密钥中 0 把改善。TYPE-I 与 NCW-shallow 重复同一模式：R2one 给出 1.45 与 1.34，R2plus 给出 0.0210 与 0.0118，94 把与 92 把中均 0 把改善。首轮减一对照保持有害符号，为 0.109 与 0.101。

扩大确认在 P2 与 P3 重复同一检验（表~\ref{tab:temporal}，图~\ref{fig:temporal}）。六个条件条目中，首轮偏移的 RR 跨 1.76--4.62（FRO-2 为 2.37--4.62），第 2 轮一次性偏移跨 1.13--1.54，持续偏移跨 0.009--0.019。没有哪个偏移处处最优：FRO-1 领跑 P2 的 RAND 列，FRO-2 在其余五列中为点估计最大值（四列单独最大，P3/TYPE-I 与 FRO-1 并列 3.27）。第 2 轮区间在六个单元中有四个单元排除 1。轮次对照变体在主分析之后加入，属于探索性结果；扩大确认的调度是预指定的。决定是否有益的是偏移的轮次与持续方式，而不是单独的幅度。

\begin{table*}[!t]
\centering
\small
\caption{P2 $(3001,37,66)$ 与 P3 $(4001,41,82)$ 上的扩大时序确认。C0 列为绝对 DFR；其余为风险比（C0 除以该调度），区间见图~\ref{fig:temporal}。每格 50 把密钥、每把 1{,}000 个错误。}
\label{tab:temporal}
\begin{tabular}{l r rrrrr}
\toprule
测试单元 & C0 DFR & FRO-1 & FRO-2 & FRO-3 & R2-2 & UNIFORM-2 \\
\midrule
P2, RAND & $2.56\times10^{-3}$ & 2.56 & 2.37 & 1.91 & 1.13 & 0.0135 \\
P2, TYPE-I & $3.14\times10^{-3}$ & 2.31 & 2.85 & 2.15 & 1.29 & 0.0188 \\
P2, NCW-shallow & $1.20\times10^{-3}$ & 1.76 & 3.00 & 2.50 & 1.36 & 0.00949 \\
P3, RAND & $1.10\times10^{-3}$ & 2.29 & 2.62 & 2.62 & 1.34 & 0.00951 \\
P3, TYPE-I & $1.70\times10^{-3}$ & 3.27 & 3.27 & 2.18 & 1.31 & 0.0165 \\
P3, NCW-shallow & $7.40\times10^{-4}$ & 3.70 & 4.62 & 3.08 & 1.54 & 0.00925 \\
\bottomrule
\end{tabular}
\end{table*}

\dblfig{fig4_temporal_dissociation.pdf}{同一偏移只在自身轮次有益，持续施加则有害。(a--c)~RAND、TYPE-I 与 NCW-shallow。上排：$P_1'$ 的主点轮次对照（灰色菱形）。下排：迁移点 P2（实心圆）与 P3（空心方块）。区间为 95\% 归档区间；虚线为风险比 1。}{fig:temporal}
\subsection{机制：不对称的假翻转抑制}\label{sec:mechanism}

第~\ref{sec:formal} 与第~\ref{sec:temporal} 小节把窗口定位在时间上；这里我们考察偏移对它触及的那一轮做了什么。一个配对探针在 RAND 下译码 20 把密钥、每把 1{,}000 个错误，得到 20{,}000 条首轮记录。平均首轮总翻转从 C0 的 29.00 降到 FRO-2 的 8.74。错误翻转下降 $9.5\times$（约为原来的十分之一），正确翻转只下降 $2.7\times$（约为原来的 38\%）。首轮后的残余错误重量从 79.02 升到 84.94。因此偏移用首轮未做完的工作，换取对假翻转的选择性。

失败走一般的译码路径。停滞占 RAND 下 C0 失败的 198/282，占 TYPE-I 的 257/372，占 NCW-shallow 的 101/141。归档分类器与预指定的四分之一预算窗口合并后没有出入，因为每次失败都用尽迭代预算。这些计数使用真实错误支持集，因此它们诊断机制，而不是给译码器一个它能自己算出的统计量。抬高阈值削掉的错误翻转远多于正确翻转；代价是正确翻转被推迟。

\dblfig{fig3_mechanism.pdf}{抬高阈值对错误翻转的抑制远强于对正确翻转。(a)~错误翻转，(b)~正确翻转，(c)~首轮后残余错误重量，单位为每次译码的平均位置数。柱为 20 把密钥、每臂 1{,}000 次译码的归档均值；每对柱上方的标签给出 FRO-2/C0。三个量使用各自独立的纵轴。}{fig:mechanism}

\subsection{迁移与稳健性}\label{sec:transfer}

把一个窗口固定在一个工作点之后，我们检验它能否经受工作点、重量、预算与语义的改变。在 $P_2=(3001,37,66)$，调度无需重新调参即完成迁移，覆盖 30 把密钥、每把 3{,}000 个错误，每变体 90{,}000 次译码。RR 在 RAND 下为 1.95 [1.58, 2.42]，TYPE-I 下为 3.21 [2.57, 4.14]，NCW-shallow 下为 2.51 [1.85, 3.46]；DFR 分别从 $2.278\times10^{-3}$ 降到 $1.167\times10^{-3}$、从 $3.000\times10^{-3}$ 降到 $9.333\times10^{-4}$、从 $1.200\times10^{-3}$ 降到 $4.778\times10^{-4}$。最差组改善为 0.384。这次运行与扩大确认是同一名义点上的两场独立战役（30 把密钥各 3{,}000 个错误，对 50 把密钥各 1{,}000 个），所以两者的 RR 不必重合。图~\ref{fig:scope}(a) 展示这三行。

次级检查保持同号。$t=94$ 与 $t=96$ 处 RR 为 2.82 [2.28, 3.60] 与 1.87 [1.72, 2.04]；最大块 $(6007,56,100)$ 给出 2.41 [2.09, 2.82]，从 $8.125\times10^{-3}$ 降到 $3.375\times10^{-3}$；预算 7、10、20 给出 1.72--3.17。12 个设置中 6 个低于可测带、6 个高于可测带。低于可测带的 5 个双方均无失败，第 6 个为 FRO-2 1/30{,}000、C0 0/30{,}000，因此低率侧无法确立方向。

高于可测带的 4 个双方均有失败且 C0 失败更多（RR 1.77、1.08、1.012 与 1.004），另外 2 个在 30{,}000/30{,}000 处饱和，没有区分力。扫描在能分辨处与该方向一致，但并不能在每一个设置上确立该方向。动态距离实现在 RAND 与 TYPE-I 下复现方向（1.59 与 2.01；$z=7.94$ 与 $11.98$），NCW 未分辨（$z=1.36$）。

最低端点把范围向下扩展：C0 在 1{,}840{,}000 次译码中失败 130 次（$7.07\times10^{-5}$），FRO-2 失败 35 次，RR 3.714 [2.69, 5.57]。在留出点 $P_4=(3607,39,76)$，一项预指定规则选中 $\delta=1$；验证记录 C0 失败 40 次，而每个一次性偏移各失败 8 次（RR 5.00）。可测设置保持方向；幅度变化，低失败率设置无法分辨。

\dblfig{fig5_scope_and_boundary.pdf}{收益在 P2 迁移，但在高交叠处消失或反向。(a)~独立 P2 迁移：C0/FRO-2 风险比与 95\% 归档区间。(b)~构造的近码字交叠 $\kappa=0,12,46,60$ 处的绝对 DFR，每处交叠 20 把密钥、每把 2{,}000 个样本（共 40{,}000 个样本）；标记的水平偏移只为避免遮挡。(c)~同一组交叠上的 C0/FRO-2 风险比；$\kappa=60$ 时两种方法均饱和，因此不画比值。}{fig:scope}
\subsection{压力构造与生产尺度证据}\label{sec:boundary}

压力实例让注入错误与近码字恰好交叠 $\kappa$ 个位置。每个构造使用 20 把密钥、每把 2{,}000 个样本，共 40{,}000 个样本。在低于块行重的交叠处偏移仍有益：对均匀构造 RR 为 3.54 [2.41, 6.00]，在 $\kappa=12$ 处为 4.75 [2.00, 19.0]。在 $\kappa=46$ 处基线崩溃：C0 在 85.31\% 的样本上失败，FRO-2 为 86.50\%，RR 0.986 [0.981, 0.991]，位于有害一侧。在 $\kappa=60$ 处两个变体对 40{,}000 个样本无一译码成功。这些绝对率以允许该构造的密钥池为条件；「捕获倾向」描述的是构造，不是已证明的捕获事件。图~\ref{fig:scope}(b,c) 汇总该边界。

近码字感知步骤把标签与结果分开。在 $\kappa=46$ 处它在 74.2\% 的样本上触发，把残余失败从 85.31\% 降到 11.07\%；在 $\kappa=60$ 处信号消失（0.76\%）。在 $P_1'$ 处该监视器完全惰性：零触发，失败计数 282 / 372 / 141 与平均迭代 4.312 / 4.511 / 4.206 均不变。因此压力族在块行重之外改变性质；浅交叠条件不构成捕获保证。

在生产块 $(r,d)=(12323,71)$ 处，朴素失败率不可触及。重量为 134 时，两个变体在 32{,}000 次译码中都记录零失败，rule-of-three 上界为 $9.4\times10^{-5}$，这是本样本的测量地板，不是最坏情况保证。只剩决策余量与收敛代价：$+2$ 把每次译码的错误翻转从 13.8 压到 2.3，同时推迟 25.5 个正确翻转；平均收敛为 2.446 对 2.054。\label{sec:production}

把重量抬到 $t^\ast=172$（比生产重量高 28\%）打开一条陡带：170 处 $1.67\times10^{-4}$，172 处 $2.17\times10^{-3}$，175 处 $1.5\times10^{-2}$，180 处 $1.39\times10^{-1}$。在该重量下窗口在 32 把密钥、每把 1{,}500 条配对错误上复现（表~\ref{tab:added}）：FRO-2 给出 RR 4.48 [2.97, 7.25]，FRONEG 为 0.162 [0.119, 0.209]。

在原生 Level-1 $(12323,71,134)$ 上，动态距离留出在偏移 0、2、4 处对两个变体均记录零失败；唯一可测信号是收敛成本。在这些偏移处，两轮收敛率在 RAND 下分别为 96.0\%、68.4\% 与 2.2\%，在 TYPE-I 下分别为 91.7\%、60.2\% 与 1.5\%，在 NCW 下分别为 93.6\%、62.3\% 与 2.0\%，每格 1{,}000{,}000 次试验。

CPA-128 预设 $(10163,71)$ 是离生产最近的一次接近。带在 $t=138$ 与 $t=142$ 之间打开；重量 140 比生产高约 4\%，C0 的 DFR 为 $8.38\times10^{-4}$；160{,}000 次配对试验中，FRO-2 给出 RR 4.96 [3.68, 7.19]，FRONEG 为 0.177 [0.148, 0.210]（表~\ref{tab:added}）。这一点是邻近性检查：它检验窗口是否在生产附近仍然存在，而不是测量 CCA 裕度。

该偏移在主工作点 $(4801,45,92)$ 的配对动态距离 C 参考中的代价很小：每变体 500{,}000 次译码、20 轮预算、三次重复下，C0 每秒 7{,}134.4 次译码、平均 4.191 次迭代，FRO-2 为 6{,}563.2 次、4.471 次（迭代比 1.067，C0/FRO-2 吞吐比 1.087），在该点约付出 7--9\% 的墙钟时间（按吞吐计 8.7\%）。公布 Level-1 常数与该点不匹配的 GRAY\_BGF 每次译码花 13.2 次迭代，并在 47.9\% 的运行上耗尽预算。这些是在 x86\_64-windows-gnu 上以 \texttt{zig cc -O2} 编译、单线程、关闭 AVX 的运行中三次墙钟时间的最小值。在已测的高交叠压力下，改善消失或反向；在生产重量下，当前样本无法解析出 DFR 收益。在生产块尺寸上，错误重量一经抬高，窗口即复现。

生产尺度的比较器配置如下。BIKE-Flip（公布常数）使用规范的 Level-1 仿射常数（$a\approx6.25\times10^{-3}$，$b\approx11.10$，$\Delta=3$；规范 v5.2 第 14 页）与本研究的 20 轮预算；BIKE-Flip（推导常数）按规范附录 A.2.4 的切线规则在测试点重新推导常数，并借用 $\Delta=3$。规范自身的预算是七轮，七轮结果与 20 轮结果并列报告。BGF（切线拟合）是 GRAY\_BGF 结构，驱动常数为测试点的 A.2.4 切线拟合仿射常数，灰步增量 $\delta=3$，20 轮预算。表~\ref{tab:added} 的所有臂都使用静态距离语义与配对实例。

\begin{table*}[!t]
\centering
\small
\caption{生产尺度检查。RR 为 C0 除以该变体，带 95\% 区间；每行给出失败数/每臂译码数。生产重量处 32{,}000 次译码零事件只给出每臂 $9.4\times10^{-5}$ 上界，不构成安全保证。}
\label{tab:added}
\begin{tabular}{llrrr}
\toprule
设置 & 变体 & 失败数 & 译码数 & RR [95\% 区间] \\
\midrule
\multirow{3}{*}{抬高重量 $(12323,71,172)$} & C0 & 94 & 48{,}000 & -- \\
 & FRO-2 & 21 & 48{,}000 & 4.48 [2.97, 7.25] \\
 & FRONEG & 582 & 48{,}000 & 0.162 [0.119, 0.209] \\
\midrule
\multirow{2}{*}{生产重量 $(12323,71,134)$} & C0 & 0 & 32{,}000 & -- \\
 & FRO-2 & 0 & 32{,}000 & -- \\
\midrule
\multirow{3}{*}{CPA-128 $(10163,71,140)$} & C0 & 134 & 160{,}000 & -- \\
 & FRO-2 & 27 & 160{,}000 & 4.96 [3.68, 7.19] \\
 & FRONEG & 756 & 160{,}000 & 0.177 [0.148, 0.210] \\
\bottomrule
\end{tabular}
\end{table*}

\subsection{现代 BF 译码器定位}\label{sec:modern}

表~\ref{tab:modern} 把窗口与更强的 BF 译码器比较。匹配条件中的失败计数为：RAND 下 282 / 68 / 83 / 287，TYPE-I 下 372 / 68 / 94 / 391，NCW-shallow 下 141 / 31 / 58 / 84，依次对应 C0、FRO-2、BIKE-Flip 与 BGF-C3。

BGF-C3 是以参数感知 CLASSIC 阈值驱动的 GRAY\_BGF 结构，灰步增量 $\delta=3$，预算 20 轮。ADMM 与再生和积译码等非 BF 译码器~\cite{watanabe2024admm} 不在本文比较范围内；它们以吞吐换取更低的 DFR，并不是普通 BF 译码器的调度变体。

BIKE-Flip 使用规范自身的阈值公式与公布的 Level-1 仿射常数；其 DFR 在每个条件中都落在 C0 与 FRO-2 之间（表~\ref{tab:modern}）。对 FRO-2 的 RAND 对比未分辨（精确 McNemar $p=0.23$）；TYPE-I 与 NCW-shallow 在 $0.0292$ 与 $0.00244$ 上给出名义分离，偏向 FRO-2。有一处不对称有利于对照方：BIKE-Flip 在本研究的 20 轮预算下运行，而规范在七轮停止。

在抬高重量处，两个 BIKE-Flip 配置分别失败 28 与 13 次，FRO-2 为 21 次（$p=0.32$ 与 $0.18$）；切线拟合 BGF 失败 63 次，C0 为 94 次（$p=0.0125$）。使用公布的 Level-1 常数时，BGF 在该重量上 94.6\% 的译码失败，这是常数错配边界。在七次迭代下，两个 BIKE-Flip 配置退化到 $1.10\times10^{-2}$ 与 $1.84\times10^{-2}$，为 C0 的六倍与九倍。

BGF-C3 对 C0 的 RR 为 0.983 [0.836, 1.156]、0.951 [0.837, 1.086] 与 1.679 [1.362, 2.111]，所以 RAND 与 TYPE-I 下与 C0 统计不可分辨，只有 NCW-shallow 给出名义分离，偏向 BGF-C3。它用 5.28、5.46、5.16 次迭代，C0 为 4.31、4.51、4.21 次，约 22\% 的迭代成本。若喂入公布的 Level-1 常数而非 $d=45$ 处推导的常数，BGF 在 38--60\% 的译码上失败，这是常数错配边界。Backflip 在其自身的动态距离语义下仍强于本文测试的任何调度~\cite{sendrier2020lowdfr,vasseur2021thesis}；这不是同语义、同调参的直接对决。在所测条件下，FRO-2 与规范译码器处于相近性能范围，但并不支配更强的 BF 译码器。这些 McNemar 检验把实例视为独立；它们是边际的，没有按密钥聚类或做多重校正。我们只把它们读作名义分离，而不是已确认的优势。

\begin{table*}[!t]
\centering
\small
\caption{$P_1'=(4801,45,92)$ 处的现代 BF 译码器比较；表中为 DFR。BGF-C3 为以参数感知 CLASSIC 阈值驱动、$\delta=3$ 的 GRAY\_BGF。$p$ 比较 BIKE-Flip 与 FRO-2（精确 McNemar）；RR 为 C0 除以 BGF-C3，带 95\% 归档区间。所有比较器共享配对静态实例；BIKE-Flip 使用 20 轮预算，而规范在七轮停止。检验为名义、实例级，未按密钥聚类或做多重校正。}
\label{tab:modern}
\begin{tabular}{lrrrrrr}
\toprule
 & \multicolumn{4}{c}{DFR} & \multicolumn{2}{c}{比较} \\
\cmidrule(lr){2-5} \cmidrule(lr){6-7}
条件 & C0 & FRO-2 & BIKE-Flip & BGF-C3 & $p$ & RR [95\% 区间] \\
\midrule
RAND & $1.022\times10^{-3}$ & $2.464\times10^{-4}$ & $3.007\times10^{-4}$ & $1.040\times10^{-3}$ & $0.228$ & 0.983 [0.836, 1.156] \\
TYPE-I & $1.319\times10^{-3}$ & $2.411\times10^{-4}$ & $3.333\times10^{-4}$ & $1.387\times10^{-3}$ & $0.0292$ & 0.951 [0.837, 1.086] \\
NCW-shallow & $5.109\times10^{-4}$ & $1.123\times10^{-4}$ & $2.101\times10^{-4}$ & $3.043\times10^{-4}$ & $0.00244$ & 1.679 [1.362, 2.111] \\
\bottomrule
\end{tabular}
\end{table*}
\subsection{为什么首轮统计量不是捷径}\label{sec:counterfactual}

轮内统计量能在产生它们的译码器内部预测失败~\cite{kim2021logistic,paiva2022pickyfix}。我们检验一个在参考变体首轮计数上训练的代理模型能否迁移到另一套调度。模型是对正确翻转数与错误翻转数做逻辑回归拟合，在 46 把密钥上训练、在 46 把不相交密钥上评估，每把 3,000 个错误。候选特征按验证集 C0 对数损失选择；预测 RR 为预测的 C0 DFR 除以预测的干预 DFR。在提供标签的臂内，它近乎完美地排序实例（AUC 0.995）；该 AUC 在用于模型选择的验证集上测得，不是最终留出测试。

跨调度时它失败了：对留出的偏移 $+1$、$+2$、$+3$，它预测 RR 0.79、0.16、0.04，而配对测量为 3.46、3.86、3.97。两个计数是处理后变量：正确翻转数偏低在 C0 下标记困难实例，在抬高阈值下却标记被推迟的翻转。因此这个特定的代理模型是产生它的调度的诊断工具，而不是另一套调度的选择器。我们不主张所有首轮统计量都失败。这正是本文直接配对测量、而不是做预测的原因。
\section{讨论与局限}\label{sec:discussion}

这些结果与一个简单的解释相一致。位置恰当的早期决策改变后续轮继承的状态。FRO-2 在严格意义上更保守：首轮翻转更少，其中错误翻转少得多，留给后续轮的真实错误更多。更轻、更干净的首轮是良性方向；更大的残差是被推迟的工作，而不是损害。因此「尽早减少错误」不是一个充分的目标；位移的符号、时点与幅度共同决定结果。

规范方的反对意见值得公平陈述。生产 BIKE-Flip 在每一轮都维持裕度，而我们的干预只触及一轮；两种主张针对不同调度。规范标定的是整条调度，而我们的破坏性情形是对逐轮标定调度施加的中途位移。FRO-2 不是规范自带的首轮保护；它是普通 BF 译码器上的一次性探针。

随机化对照通过把位置从调度中移除来补全论证。这一对照是描述性的：运行在 BF30（筛选点的固定阈值 BF 变体）上，并在每一轮为阈值加入掷币式抖动。对称抖动把失败率从 4.29\% 推到 19.67\%；仅正向的单侧配方落在 50.5\% 与 51.8\% 之间。同一配方用于普通 BF 译码器时，小样本探针未留下可分辨差异（Fisher $p\ge0.54$）；我们把它读作该探针内未见分辨效应，而不是随机化普遍无效的证据。无结构混合不能替代有定位的确定性调度。

本研究有四个边界。全部主要结论依托约 $10^{-4}$ 到 $10^{-3}$ 的失败率，因此不外推生产目标。两个工作点承载主要推断，第三个为补充；其它块尺寸只有行为筛查。偏移网格与被筛调度是我们的构造，跨实现检查是同项目而非独立复现。诊断标签使用真实错误支持集，部署者并不知道它。在生产参数下，抬高重量的窗口比生产重量 134 高 28\%；生产重量证据是决策余量与收敛成本，零失败上界为 $9.4\times10^{-5}$。捕获未被证明；交叠是压力构造。未来工作应把测量推向生产尾部，并在更大密钥池上测试独立实现。

\section{结论}\label{sec:conclusion}

普通 BF 译码器的首轮是一扇陡峭且不对称的双向窗口。保守偏移有益，而宽松侧在两个计数内崩塌。决定条件失败风险的是偏移的轮次与持续方式，而不是单独的幅度。窗口与不对称的假翻转抑制相一致。干预剔除错误决策，而不是加速纠错。它无需重调即跨工作点与重量迁移，在所测捕获倾向构造处止步，并在生产块尺寸上于错误重量抬升后复现；在生产重量本身，样本只给出 rule-of-three 上界、决策余量与收敛代价。一个首轮代理模型未能预测另一套调度的效应，随机化也不能复现有定位决策的收益。实践上，逐变体配对实例、以密钥为聚类单元做推断，并在失败可测处标定阈值。
% ===========================================================================
\section*{数据与代码可用性}
\addcontentsline{toc}{section}{数据与代码可用性}
实现、冻结配置、种子派生规则、逐密钥原始计数与分析代码归档于 \url{https://doi.org/10.5281/zenodo.23636424}，开发镜像见 \url{https://github.com/Leekj10/qc-mdpc-experiment}。图源数据表与把正文引用的每个数值与归档计数逐一重新核对的校验脚本随同一归档发布。

\bibliographystyle{IEEEtran}
\bibliography{refs}

\end{document}
